\documentclass[11pt]{article}
\usepackage[margin=0.85in]{geometry}
\usepackage{amsmath,amssymb,mathtools}
\usepackage{booktabs,array,enumitem,microtype}
\setlength{\parindent}{0pt}
\setlength{\parskip}{0.65em}

\newcommand{\ket}[1]{\lvert #1\rangle}
\newcommand{\bra}[1]{\langle #1\rvert}
\newcommand{\I}{\mathbb{I}}
\newcommand{\CNOT}[2]{\operatorname{CNOT}_{#1\to #2}}

\begin{document}

\begin{center}
  {\Large\bfseries Problem 1 -- Question 6}\\[0.3em]
  {\large Measurement}
\end{center}

\section*{6(a): Joint measurement of \(Z_1Z_2\) and \(X_1X_2\)}

The two observables commute:
\[
 (Z_1Z_2)(X_1X_2)
 =(Z_1X_1)(Z_2X_2)
 =(X_1Z_1)(X_2Z_2)
 =(X_1X_2)(Z_1Z_2).
\]
They can therefore be measured jointly without an uncertainty
conflict.

\paragraph{\(Z_1Z_2\) syndrome.}
Prepare an ancilla \(a_Z\) in \(\ket0\), apply
\[
 \CNOT{q_1}{a_Z},\qquad \CNOT{q_2}{a_Z},
\]
and measure \(a_Z\) in the \(Z\) basis. If the measured bit is \(m_Z\),
the measured eigenvalue is
\[
 z=(-1)^{m_Z}.
\]
The data state is projected by
\[
 \Pi^{(ZZ)}_{z}=\frac{\I+zZ_1Z_2}{2}.
\]

\paragraph{\(X_1X_2\) syndrome.}
Prepare a second ancilla \(a_X\) in \(\ket+\) by applying \(H\) to
\(\ket0\). Apply
\[
 \CNOT{a_X}{q_1},\qquad \CNOT{a_X}{q_2},
\]
then apply \(H\) to \(a_X\) and measure it in the \(Z\) basis. If the
result is \(m_X\), the measured eigenvalue is
\[
 x=(-1)^{m_X},
\]
and the data state is projected by
\[
 \Pi^{(XX)}_{x}=\frac{\I+xX_1X_2}{2}.
\]
The two syndrome extractions may be done sequentially with one resettable
ancilla, or with two ancillas.

Conditioned on outcomes \(m_Z=m_X=0\), the normalized output is
\[
\boxed{
\ket{\psi_{++}}
=\frac{(\I+Z_1Z_2)(\I+X_1X_2)\ket\psi}
{\left\|(\I+Z_1Z_2)(\I+X_1X_2)\ket\psi\right\|}.
}
\]
This is the joint \(+1\) projection. The success probability is
\[
 p_{++}
 =\bra\psi
 \frac{(\I+Z_1Z_2)(\I+X_1X_2)}{4}
 \ket\psi.
\]
For two data qubits, the joint \(+1\) eigenspace is one-dimensional and
is spanned by
\[
 \ket{\Phi^+}=\frac{\ket{00}+\ket{11}}{\sqrt2}.
\]

If deterministic preparation of \(\ket{\Phi^+}\), rather than literal
postselection, is desired, use classical feedback: apply \(X_1\) when
\(m_Z=1\), and apply \(Z_1\) when \(m_X=1\). \(X_1\) flips only the
\(ZZ\) eigenvalue, while \(Z_1\) flips only the \(XX\) eigenvalue.

\section*{6(b): Projective measurement of a general Pauli operator}

Let
\[
 P=\eta\,P_1\otimes\cdots\otimes P_n,\qquad
 \eta\in\{+1,-1\},\qquad
 P_j\in\{\I,X,Y,Z\}.
\]
Ignore identity positions. On every nonidentity position, rotate its
Pauli operator to \(Z\):
\[
 B_j=
 \begin{cases}
 H,&P_j=X,\\
 HS^\dagger,&P_j=Y,\\
 \I,&P_j=Z.
 \end{cases}
\qquad
B_jP_jB_j^\dagger=Z.
\]
Chronologically, \(Y\)-basis conversion applies \(S^\dagger\) and then
\(H\).

Prepare an ancilla \(a\) in \(\ket0\). After the basis changes, apply a
CNOT from every nonidentity data qubit to \(a\), then measure \(a\) in
the \(Z\) basis. The measured bit \(m\) gives eigenvalue
\[
 p=\eta(-1)^m.
\]
Finally undo the basis changes if the data must remain available.
The corresponding projectors are
\[
\boxed{\Pi_{\pm}=\frac{\I\pm P}{2}.}
\]
This is a nondestructive projective measurement: within a fixed
eigenspace, superpositions are preserved. It uses one ancilla and
\(\operatorname{wt}(P)\) CNOTs. If destroying the data is acceptable,
one may instead measure every nonidentity qubit in its local Pauli basis
and multiply the individual eigenvalues.

\section*{6(c): Reducing the GHZ-like state to one qubit}

Write
\[
 \ket{\Psi_n}
 =\alpha\ket{0}^{\otimes n}
 +\sqrt{1-\alpha^2}\ket{1}^{\otimes n}.
\]
Keep \(q_1\). Measure \(q_2,\ldots,q_n\) in the \(X\) basis by applying
Hadamards and then measuring in the computational basis. Let the outcomes
be \(m_2,\ldots,m_n\), and define their parity
\[
 p=m_2\oplus\cdots\oplus m_n.
\]
Because
\[
 \bra{\pm}0\rangle=\frac1{\sqrt2},\qquad
 \bra{\pm}1\rangle=\frac{\pm1}{\sqrt2},
\]
the unmeasured qubit is left in
\[
 \alpha\ket0+(-1)^p\sqrt{1-\alpha^2}\ket1.
\]
Apply the classically controlled correction \(Z^p\) to \(q_1\). The
final state is
\[
\boxed{
\alpha\ket0+\sqrt{1-\alpha^2}\ket1.
}
\]
Every outcome string is acceptable; the parity correction removes its
known phase. Thus the protocol succeeds deterministically and requires
no ancillary qubits.

\end{document}
